% Source: pca_and_reconstruction.tex
\begin{notechapter}
\chapter{PCA and Data Reconstruction}\label{ch:pca-reconstruction}

\section{Setup}
Let \( X \) be a \textbf{centered} \( n \times m \) data matrix where:
\begin{itemize}
  \item Each row represents a data point in \( \mathbb{R}^m \).
  \item The mean of each column is \( 0 \).
\end{itemize}

\subsection{Covariance Matrix}
The covariance matrix is \( \frac{1}{n}X^T X \) (not \( XX^T \)),
which is \( m \times m \):
\[
  \text{Cov}(X) = \frac{1}{n}X^T X
\]
Its eigenvectors define the principal components of \( X \).

\subsection{Eigenvectors}
Let \( V \) be the \( m \times m \) orthogonal matrix whose columns
are the eigenvectors of \( \frac{1}{n}X^T X \). These eigenvectors:
\begin{itemize}
  \item Are orthonormal: \( VV^T = V^T V = I \).
  \item Form a basis for \( \mathbb{R}^m \).
\end{itemize}

\section{Projection and Reconstruction}
\subsection{Projection}
The coordinates of the data in the eigenbasis are:
\[
  U = X V
\]
Each row of \( U \) is a data point in the eigenvector basis.

\subsection{Reconstruction}
The original data can be reconstructed using:
\[
  X = U V^T
\]
\textbf{Proof:}
\[
  U V^T = (X V)V^T = X (V V^T) = X I = X
\]

\section{2D Example}
\subsection{Data Matrix}
Let \( X \) be a centered \( 3 \times 2 \) matrix:
\[
  X = \begin{bmatrix}
    -1 & -1 \\
    0 & 0 \\
    1 & 1
  \end{bmatrix}
\]

\subsection{Covariance Matrix}
Compute \( X^T X \):
\[
  X^T X = \begin{bmatrix}
    2 & 2 \\
    2 & 2
  \end{bmatrix}
\]

\subsection{Eigenvectors}
Solve \( \det(X^T X - \lambda I) = 0 \):
\[
  \det\begin{bmatrix}
    2 - \lambda & 2 \\
    2 & 2 - \lambda
  \end{bmatrix} = \lambda^2 - 4\lambda = 0 \implies \lambda_1 = 4, \, \lambda_2 = 0
\]
Eigenvectors:
\[
  \bm{v}_1 = \frac{1}{\sqrt{2}}\begin{bmatrix} 1 \\ 1 \end{bmatrix},
  \quad
  \bm{v}_2 = \frac{1}{\sqrt{2}}\begin{bmatrix} 1 \\ -1 \end{bmatrix}
\]
Thus:
\[
  V = \frac{1}{\sqrt{2}}\begin{bmatrix}
    1 & 1 \\
    1 & -1
  \end{bmatrix}
\]

\subsection{Projection}
Compute \( U = XV \):
\[
  U = \frac{1}{\sqrt{2}}\begin{bmatrix}
    -2 & 0 \\
    0 & 0 \\
    2 & 0
  \end{bmatrix}
\]

\subsection{Reconstruction}
Verify \( UV^T = X \):
\[
  UV^T = \frac{1}{\sqrt{2}}\begin{bmatrix}
    -2 & 0 \\
    0 & 0 \\
    2 & 0
  \end{bmatrix} \cdot \frac{1}{\sqrt{2}}\begin{bmatrix}
    1 & 1 \\
    1 & -1
  \end{bmatrix}^T = X
\]
(The eigenvalues of the covariance matrix $\frac{1}{3}X^TX$ are $4/3$ and $0$: the same eigenvectors, eigenvalues divided by $n$.)

\section{Truncated Reconstruction}
The interest of PCA is that one can keep only the first $k < m$ components. Let $V_k$ be the
$m\times k$ matrix of the eigenvectors associated with the $k$ largest eigenvalues
$\lambda_1 \geq \dots \geq \lambda_k$ of $\frac{1}{n}X^TX$, and $U_k = XV_k$ ($n\times k$). The
rank-$k$ reconstruction is
\[
  \hat{X}_k = U_k V_k^T = X V_k V_k^T ,
\]
i.e.\ each data point is orthogonally projected on the span of the first $k$ principal
directions ($V_kV_k^T$ is the projection matrix). Its mean squared reconstruction error is
\[
  \frac{1}{n}\big\|X - \hat{X}_k\big\|_F^2 = \sum_{j=k+1}^{m} \lambda_j ,
\]
the sum of the discarded eigenvalues, and no rank-$k$ matrix does better (Eckart--Young theorem).

In the example, $\lambda_2 = 0$: keeping only $\bm{v}_1$ already reconstructs $X$ exactly,
\[
  \hat{X}_1 = X\bm{v}_1\bm{v}_1^T = \frac{1}{\sqrt{2}}\begin{bmatrix} -2 \\ 0 \\ 2 \end{bmatrix} \frac{1}{\sqrt{2}}\begin{bmatrix} 1 & 1 \end{bmatrix} = X ,
\]
because all the points lie on the line spanned by $\bm{v}_1$.

\section{Conclusion}
If \( V \) is the eigenvector matrix of \( \frac{1}{n}X^T X \), then:
\[
    \boxed{UV^T = X}
\]
and keeping only the first $k$ eigenvectors gives the best rank-$k$ approximation $XV_kV_k^T$, whose error is the sum of the discarded eigenvalues.
\end{notechapter}
